\documentclass[UTF8,a4paper,12pt]{ctexart}

\usepackage{amsmath}
\usepackage{amssymb}
\usepackage{geometry}
\usepackage{hyperref}
\usepackage{enumitem}
\usepackage{xcolor}
\usepackage{listings}

\geometry{left=2.5cm,right=2.5cm,top=2.5cm,bottom=2.5cm}
\hypersetup{
    colorlinks=true,
    linkcolor=black,
    urlcolor=blue
}
\setlist{nosep}

\lstset{
    basicstyle=\ttfamily\small,
    breaklines=true,
    frame=single,
    columns=fullflexible
}

\title{Milestone 报告：基于 SDF 隐空间扩散模型的三维形状生成}
\author{符海宝\hspace{1cm}姜书伟\hspace{1cm}刘绍凡}
\date{\today}

\begin{document}

\maketitle

\section{选题背景}

三维内容生成是几何计算与生成式模型交叉中的重要问题，目标是从随机噪声、类别条件、文本、图像或部分几何观测中生成可用的三维形状。相比二维图像生成，三维生成不仅需要保证形状的视觉合理性，还需要满足空间结构一致性、可渲染性以及后续编辑、重建和仿真的需求。

本项目计划研究一种基于 SDF latent diffusion 的三维形状生成方法。与直接在点云或高分辨率 SDF 网格上训练扩散模型相比，隐空间扩散先通过三维自编码模型将形状压缩到较低维的 latent space，再在该隐空间中训练扩散模型，可以降低计算开销并提升生成过程的稳定性。该路线与 SDFusion 等工作思路接近，但本阶段不直接复现完整 VQ-VAE 版本，而是先实现一个更轻量、便于验证的 3D VAE + latent DDPM 基础通路。

本项目的基础目标是实现一个单类别、无条件的三维形状生成流程：从 ShapeNet chair 类别 mesh 出发，将其转换为 T-SDF grid，训练可解码的三维隐空间表示，并在 latent space 中训练 DDPM，最终从随机噪声生成 SDF / mesh。当前阶段为了优先验证代码接口与方法通路，先使用 synthetic SDF 数据完成端到端 smoke test，真实 ShapeNet 数据接入作为下一步工作。

\section{核心方法}

本项目采用 ``SDF 表达 + 隐空间扩散'' 的路线。首先将三维形状表示为规则网格上的 SDF 或 truncated SDF。相较于点云，SDF 更适合表达封闭表面，也可以通过 Marching Cubes 直接提取 mesh；相较于直接在高分辨率 SDF 上扩散，先压缩到隐空间可以降低计算量。

在隐空间建模阶段，项目使用一个轻量 3D VAE。Encoder 由多层 3D convolution block 组成，将 SDF grid 压缩为较小的 latent grid；decoder 则通过 3D transposed convolution 将 latent 还原为 SDF。VAE 的训练目标以 SDF 重建误差为主，并加入较小权重的 KL loss，使 latent 具有一定连续性和可采样性。

在生成阶段，扩散模型不直接处理原始 SDF，而是在 VAE 得到的 latent grid 上训练。模型先对 clean latent 加入高斯噪声，再使用带时间步条件的 3D U-Net 预测噪声并学习反向去噪过程。选择 3D U-Net 而不是 Transformer 的原因是，本项目的 latent 仍然是规则三维网格，3D U-Net 更直接、轻量，适合作为课程项目的基础实现。

采样时，模型从随机高斯 latent 出发，经过 DDPM 反向去噪得到生成 latent，再由 VAE decoder 解码成 SDF，最后使用 Marching Cubes 提取三维 mesh。

整体通路如下：
\begin{lstlisting}
SDF Grid X
  -> 3D VAE Encoder
  -> Latent Grid z0
  -> Forward Diffusion: z0 -> zt
  -> 3D U-Net Denoiser epsilon_theta(zt, t)
  -> Reverse Diffusion: zT -> z0
  -> VAE Decoder
  -> Generated SDF
  -> Marching Cubes
  -> Generated Mesh
\end{lstlisting}

\section{可拓展方法}

后续可在当前隐空间扩散框架上继续扩展以下方向：
\begin{enumerate}
    \item 将普通 3D VAE 扩展为 VQ-VAE 或 codebook 版本，使方法更接近 SDFusion 一类模型。
    \item 使用 DDIM 采样替代标准 DDPM 采样，减少采样步数并提升推理效率。
    \item 加入类别条件，将 chair、airplane、car 等 ShapeNet 类别通过 class embedding 注入 3D U-Net，实现多类别条件生成。
    \item 引入文本、图像或 partial shape 条件，通过 CLIP、CNN 或 shape encoder 提取条件特征，并注入去噪网络。
    \item 在生成结果稳定后，对不同 SDF 分辨率、latent 分辨率和采样步数进行对比实验。
\end{enumerate}

\section{作业进展}

目前已经完成基础代码框架，并跑通了最小端到端方法通路。当前实现使用 synthetic SDF 数据进行 smoke test，验证了如下流程：
\begin{lstlisting}
synthetic SDF -> 3D VAE -> latent DDPM
-> 3D U-Net denoiser -> VAE decoder
-> Marching Cubes mesh export
\end{lstlisting}

已经完成的主要模块包括：
\begin{itemize}
    \item synthetic SDF 数据生成、3D VAE、3D U-Net、GaussianDiffusion3D 和 Marching Cubes 导出工具。
    \item \texttt{smoke\_test.py}、\texttt{train\_vae.py}、\texttt{train\_diffusion.py} 和 \texttt{sample.py} 等基础脚本入口。
    \item smoke test 已能够生成输入、重建和采样得到的 \texttt{.ply} mesh 文件，并保存基础 checkpoint。
\end{itemize}

当前 smoke test 运行命令为：
\begin{lstlisting}
python scripts/smoke_test.py --config configs/smoke.yaml
\end{lstlisting}

下一步 TODO 如下：
\begin{enumerate}
    \item 完成 ShapeNet chair 数据接入，实现真实 mesh 的读取、归一化和 T-SDF grid 预处理。
    \item 在真实或更丰富的 SDF 数据上训练 VAE，提升 reconstruction mesh 质量。
    \item 批量提取并缓存 VAE latent，用于更稳定地训练 latent diffusion。
    \item 增加 diffusion 训练步数，观察生成 mesh 是否逐步形成合理 chair 结构。
    \item 补充可视化脚本，展示 input、reconstruction、generated mesh 的对比图。
    \item 在基础结果稳定后，再尝试 DDIM、类别条件生成或 VQ-VAE 扩展。
\end{enumerate}

\end{document}
